% =============================================================================
%  一、问题陈述
% =============================================================================
\section{问题陈述}

\subsection{研究背景：从“投入多少”到“效率多高”}

高等教育的评价对象正在从“教得好不好”扩展到“学得会不会”。随着学分制与自
主学习的推行，同一门课程里学生可自由支配的课外时间显著增加，学习结果却呈
现出很大的离散性：有人每天投入四五个小时却收效有限，有人在较短时间内完成
听说读写的闭环并取得高分。教育心理学把这种差异归因于学习策略，即学生在时
间分配、复习方式、应试训练、心理调适与兴趣拓展上的组合选择。于是，教学管
理者真正关心的问题并不是“谁更努力”，而是“在给定的投入组合下，谁的转化效
率更高；效率低的人应当朝哪个方向、以多大强度调整，才能以可接受的时间代价
上一个台阶”。

传统上回答这一问题依靠三类工具：其一是相关分析与回归，它给出“某行为与成
绩同向变动”的平均结论，却无法处理多投入、多产出之间的替代关系；其二是问
卷描述统计，它告诉我们学生“做了什么”，不回答“做得值不值”；其三是聚类或
潜类分析，它把学生分成若干类型，但类型本身没有优劣，无法直接转化为建议。
数据包络分析（Data Envelopment Analysis, DEA）\refs{1} 的引入改变了这一点：
它把每个学生视为一个决策单元（Decision Making Unit, DMU），用线性规划在
“投入—产出”空间中构造前沿面，把个体的相对效率写成投入可压缩的比例，从而
在同一标尺下比较“转化能力”。这一思路在教育效率、学校绩效与资源配置中被广
泛使用。

然而，标准 DEA 只回答“你现在离前沿有多远”，不回答“你离前沿远是因为起点
差，还是因为路径差”。为区分这两种情形，Seiford 与 Zhu 提出
context-dependent DEA\refs{2}：先剔除强有效单元，再在剩余集合上重新评价，
如此逐层剥离，得到一个由外向内的效率分层结构，并用吸引力、进步度、邻近度
三个指标描述个体在层间的相对位置。这一想法被后续研究用于技术效率的分级、
区域经济梯度与医院绩效的阶梯式改进。

\subsection{被解析的论文及其贡献}

本文解析的对象是 Zhao 等人 2023 年发表于 \textit{Expert Systems with
Applications} 第 211 卷的研究\refs{3}。该文以北京某高校 519 名非英语专业本
科生的英语课程为背景，围绕 14 个学习行为题项（兴趣、复习、考试次数、临场
状态、心理、休息、模拟考、背诵、听力、口语、阅读、课外拓展、辅导、出国意
愿）构造了一个名为 CHPI 的推荐框架，其流程可概括为四个步骤：

\begin{enumerate}[label=（\arabic*）,itemsep=2pt,topsep=2pt,leftmargin=3em]
\item \textbf{复合量化}：不同题项的量纲与题型并不一致（五级有序、计数、连
  续时长），原文用“累积百分比法”把每一题的作答映射为 0—100 的得分，再按题
  型加权合成 4 个投入指标与 1 个产出指标，并以题项方差 $\sigma^2<10$ 作为
  “信息量不足需重赋权”的判据；
\item \textbf{本体聚类}：以 Gower 距离处理混合数据类型，用 PAM
  （$k$-medoids）把学生分为 5 类，理由是“不同学习风格的学生不应放在同一个
  前沿面上比较”；
\item \textbf{类内效率评价}：在每一类内部用投入导向的 CCR 模型计算效率
  $\theta_i$，并用 $\varepsilon$ 松弛与 Archimedes 无穷小处理弱有效；
\item \textbf{分层与路径推荐}：用 context-dependent DEA 的剔除式算法把每类
  再分为 5 层，计算吸引力 $\Omega$、进步度 $P$、邻近度 $D$ 与障碍度 $H$，
  再按“上一层中 $H$ 最小、距离最近”的规则为个体拼接一条逐层上升的参考路
  径。原文以 5 类中的第 5 类、效率 0.6771 的学生 S480 为例，给出
  S480$\to$S502$\to$S216$\to$S18 的四步路径，单步提升 25.38\%，累计提升
  43.97\%。
\end{enumerate}

这篇论文的价值在于把“效率评价”推进到“效率分层”，再把“分层”推进到“个体
路径”，在方法论上完成了从静态打分到结构化建议的跨越；在应用上，它给出了
一套可执行的数据处理口径（尤其是混合题型的统一量化），这对教育数据挖掘领
域具有直接的借鉴意义。也正因为它在这条链条上走得足够远，链条上每一环的薄
弱之处才暴露得格外清楚。

\subsection{原文的三处结构性局限}

\textbf{第一，评价是静态的，缺少时间维度。}原文的效率 $\theta_i$ 是某一学
期截面数据的比值，层级是个体在这一截面中的相对位置。于是它只能回答“谁在
这一层”，不能回答“从这一层到上一层平均需要多久”“这一跃迁的概率有多大”
“跳上去之后能否稳住”。对教学干预而言，这三个问题恰恰是最要紧的：一个需要
八周、成功率六成的建议，与一个需要一年、成功率一成的建议，其优先级完全不
同。

\textbf{第二，层级依赖参照集，不可跨群体迁移。}剔除式分层把“被当前前沿
支配的单元”留在下一层，层数由数据本身的效率分布决定，而不是由某个外生标
准决定。这意味着同一名学生在不同班级、不同学期、不同题项权重下会被分到不
同的层；层级不是一把可以随身携带的尺子，而是一张只对当前样本有效的排名
表。原文把层号直接写进推荐语（“从第 5 层升到第 4 层”），在跨群体使用时缺
少可解释性。

\textbf{第三，推荐机制本质上是“模仿”，而非“最优控制”。}路径规则是在上一
层中找 $H$ 最小、$L_1$ 距离最近的同伴，令个体向其行为组合靠拢。这一规则隐
含两个强假设：其一，同样的行为组合在不同个体身上产生同样的效率（忽略个体
异质性与边际收益递减）；其二，越接近高效同伴就越优（忽略了“投入已经饱和、
再加投入反而稀释效率”的可能）。更实际的问题是，它只告诉学生“跟谁学”，不
告诉学生“学多少、练多久、先做哪一项”。

除这三处结构性局限之外，本文在核对原文表格时还发现一处内部不一致：原文
文字部分把障碍度定义为 $H=1/\Omega+P$，但表 3 中给出的 $H$ 值与同一行的
$(\Omega,P)$ 无法对应。例如某个体 $\Omega=0.0514$、$P=0.0096$，按定义应得
$H\approx 19.5$，而表中 $H=0.43$。这说明原文在指标合成或表格转录环节存在未
披露的处理，也说明“三准则”的可复现性需要逐项检验，而不能照抄公式。

\subsection{本文的任务与组织结构}

基于上述判断，本文把工作拆成三个前后衔接、权重相当的部分，并以“能否被独立
复算”作为唯一的验收标准。

\begin{enumerate}[label=（\arabic*）,itemsep=3pt,topsep=3pt,leftmargin=3em]
\item \textbf{破解}：把原文从“问卷题项”到“推荐路径”的全过程拆成“输入—变
  换—输出”的证据链，逐步写出数学形式、算法伪码与隐含假设，标注每一步的
  自由度（哪些参数由数据决定、哪些由作者选择、哪些未披露），从而回答“这篇
  论文究竟做了什么、凭什么成立”。
\item \textbf{复现}：在拿不到原始问卷的前提下，以题项分布、样本规模、类规
  模、以及原文公开的锚点（S480 的效率 0.6771、单步 25.38\%、累计
  43.97\%）为约束，逆向重建一份可计算的数据集，用 R 语言实现
  PAM-Gower 聚类、类内 CCR-DEA、分层与三准则，逐项与原文对照，并对偏差给
  出机制性解释。
\item \textbf{替换}：不再使用 context-dependent DEA，而用随机过程重建同一
  主题。把五级效率层级视为状态空间，建立带协变量的生灭过程，估计转移强度，
  解析求解首达时间，用 Gillespie 算法校验；再用 Ornstein-Uhlenbeck 扩散刻
  画连续效率的随机演化，最后在学时预算约束下求解“哪几项动作、各投入多少”
  的最优策略。与原文的静态排序相比，这套方法输出的是“概率—时间—策略”的
  三元组。
\end{enumerate}

三部分的技术路线与交付物见表~\ref{tab:route}。需要说明的是，第一、二部分
并非“铺垫”，它们承担两项不可替代的功能：其一，只有把原文的每一步都复算出
来，替换方法时才能保证“数据、主题、评价对象”三者不变，从而使新旧方法的差
异可归因于方法本身；其二，复现过程中暴露的口径问题（层级定义、指标合成、
聚类数选择）正是随机过程建模必须重新回答的问题——例如“状态”究竟应当是剔除
式层级，还是可以保证占满状态空间的分位带层级。

\begin{table}[htbp]
\centering
\tabfont
\caption{本文三阶段任务、方法与交付物}
\label{tab:route}
\begin{tabularx}{\textwidth}{@{}p{1.5cm}p{3.7cm}Xp{3.6cm}@{}}
\toprule
阶段 & 输入 & 主要方法 & 交付物 \\
\midrule
破解 & 原文 PDF（12 页）、
公开表格 & 证据链拆解、公式重述、算法还原、自由度标注 &
  流程图、公式表、可复现性判定表 \\
复现 & 题项分布与锚点约束 &
  累积百分比量化、Gower 距离、PAM、CCR-DEA、分位带与剔除式分层 &
  519 人数据集、效率与层级、锚点对照表、9 幅插图 \\
替换 & 同一数据集、同一主题 &
  生灭过程、极大似然、鞅停时、首达时间、Poisson 一致化、反射原理、Itô 公式、OU 扩散、预算约束优化 &
  转移强度、时间尺度、跃层概率、最优策略、新旧对比表 \\
\bottomrule
\end{tabularx}
\end{table}

全文余下部分安排如下：第二节把实际问题翻译为数学问题，并说明为什么随机过
程是合适的语言；第三节是主体，依次给出原文解析、CHPI 复现、随机过程重构与
系统对比；第四节总结结论、局限与可延伸的方向。

